% !PRINTED 671
\weberpage{671}
%<<CONT
entsprechender Weise die Punkte aufnehmen, in denen $\eta$ eine negative Ordnungszahl hat. Dann sind die Polygone $\mathfrak{A}$ und $\mathfrak{B}$ von gleicher Ordnung, nämlich von der Ordnung der Funktion $\eta$, und durch diese Polygone ist die Funktion $\eta$ bis auf einen konstanten Faktor bestimmt.

Wir können daher die symbolische Bezeichnung einführen:

\weq{1}{\eta = \dfrac{\mathfrak{A}}{\mathfrak{B}}\,,}

und $\mathfrak{A}$ den \gesperrt{Zähler} oder das \gesperrt{Obereck}, $\mathfrak{B}$ den \gesperrt{Nenner} oder das \gesperrt{Untereck} von $\eta$ nennen. Nach dieser Definition sind zunächst $\mathfrak{A}$ und $\mathfrak{B}$ relativ prim zueinander. Wir wollen aber diese Bezeichnung noch dadurch erweitern, daß wir, wenn $\mathfrak{M}$ ein beliebiges Polygon ist:

\weq{2}{\dfrac{\mathfrak{M}\mathfrak{A}}{\mathfrak{M}\mathfrak{B}} = \dfrac{\mathfrak{A}}{\mathfrak{B}}}

setzen. Dann ist die Bezeichnung (1) von der Beschränkung frei, daß $\mathfrak{A}$ und $\mathfrak{B}$ relativ prim sein sollen. Beide Polygone sind immer noch von gleicher Ordnung, aber diese kann größer sein als die Ordnung von $\eta$.

Bei dieser Darstellung der Funktionen $\eta$ gelten dann für die Multiplikation und Division dieselben Regeln, wie beim Rechnen mit Zahlenbrüchen im Gebiete der natürlichen Zahlen.

In (1) können $\mathfrak{A}, \mathfrak{B}$ nicht beliebige $m$-Ecke sein, und die Erforschung der Beziehung zwischen diesen ist die große Frage, die, in anderer Weise, durch das \gesperrt{Abelsche Theorem} beantwortet wird.

Wir stellen folgende Definition auf:

1. \gesperrt{Können zwei $n$-Ecke $\mathfrak{A}, \mathfrak{A}'$ Obereck und Untereck einer Funktion $\eta$ in $\Omega$ sein, so heißen $\mathfrak{A}$ und $\mathfrak{A}'$ äquivalent.}

In Zeichen: Gibt es eine Funktion $\eta$, der die Bezeichnung

\[
\eta = \dfrac{\mathfrak{A}}{\mathfrak{A}'}
\]

zukommt, so ist

\weq{3}{\mathfrak{A} \sim \mathfrak{A}'.}

Aus der Formel

\[
\eta = \dfrac{\mathfrak{A}}{\mathfrak{A}'}\,, \quad \eta' = \dfrac{\mathfrak{A}}{\mathfrak{A}''}\,, \quad \dfrac{\eta'}{\eta} = \dfrac{\mathfrak{A}'}{\mathfrak{A}''}
\]

ergibt sich der Satz:
%CONT>>
